\subsection{幼主与守成的困难}

后周末年的幼主继位，使陈桥发生的背景更易理解。守成并不是把已建立的制度放在那里不动；它需要让军队相信供给与指挥仍有效，让官员相信任命与法令仍有连续性，让地方相信中枢不会因继承而无力保护道路和市场。若这些信念同时变弱，某一支控制近畿兵权的集团便可能把“保全秩序”的语言转化为接管权力的理由。

960年的陈桥不必被写成预定的天命，也不必缩减为一次孤立兵变。它发生在后周刚经历北伐、南征和整军的条件下，军政资源集中于开封周边，继承又使这些资源的归属不够确定。赵匡胤建立宋，见\eventref{LT-960-118}{陈桥与宋代建立（960）}，使中原中枢得到新的持有者；新政权仍须向旧有官员、军队和地方说明，哪些关系将被保留、哪些负担将被重新分配。政权转换的成功，包含了让多数日常事务继续运转的能力。

\section{战后接收：占一座城与用一座城是两回事}

923年后唐入汴、925年入蜀，两次行动都可被写成“灭国”，但它们之后的工作并不相同。汴州是后梁都城，官署、仓储、市场和旧军集中，胜者进入后须马上决定谁保管府库、谁看守城门、谁处理旧臣和居民。成都则处在另一套平原与山路的网络中，后唐军队能抵达城下，不代表它已掌握周边交通、盐货和地方军政。两个“灭”字遮蔽了不同的接收成本。

战后最先需要恢复的常常不是宏大礼仪，而是可预期性。市民需要知道市场是否开放，商旅需要知道关津由谁控制，地方官需要知道文书交给谁，士卒需要知道粮食从哪里来。新朝若对旧有人员一概清除，可能失去了解账目、道路和地方关系的人；若一概保留，又要承担旧网络重新结盟的风险。接收因此是一种反复权衡，不会因在宫殿中举行受降而结束。

南唐进入闽地和湖南也面临类似而不相同的问题。945年建州归入南唐，952年南唐介入楚地，见\eventref{LT-945-104}{南唐取闽（945）}与\eventref{LT-952-SOUTHERN-TANG-CHU}{南唐介入湖南（952）}。福建的山路、港口和地方军政使接收难以照搬江淮经验；湖南的水系、军镇与潭州中心又不同于建州。扩张让南唐的名义与军队到达新的区域，也增加了驻防、运输、任命和安抚的需求。强大不在于能否一时攻下，而在于能否让这种需求不压垮已有的资源。

后周接收南唐江北诸州时，问题再次出现。957年的称臣与割地给后周提供了新的边界和城市，却没有把地方账目、市场关系或守军忠诚整齐地移交。后周须判断哪些旧官吏可以使用，驻军从何得到供应，河道和渡口怎样与中原网络相连；南唐则要在失地后重建防线和信誉。条约的文字可以固定两国之间的关系，地方的接续却只能在时间中慢慢完成。

\section{工匠、船户与书手：国家能力中的非官员}

五代十国的叙事常由帝王、将领和大臣推动，但国家的可见能力常由不在本纪中心的人完成。修城的工匠知道墙体和材料，疏浚河道的人知道水势，船户知道浅深与风向，书手和胥吏知道如何让一份命令成为可执行的记录，市场中的经纪人与商人维持货物和信用的流转。他们未必拥有决定国策的权力，却使国策能否越过宫门成为问题。

例如，船户在江淮、两浙、岭南和长江中游拥有不同的知识。运河上的船运受堤闸、河道和官府调度影响，钱塘江与沿海航行还要面对潮汐和风季，珠江与交州河口则需要熟悉支流和浅滩。国家可以征用船只或规定航道，却要承担由此造成的延误、逃避或市场断裂。我们很少有船户留下的完整账簿，不能为他们虚构一致的立场；但任何把水路写成无人的箭头的叙事，都漏掉了这一层。

书手与胥吏的连续性同样重要。一个政权可以迅速改元、换官、铸钱，却需要有人持续保管户籍、文案、仓账和契约。五代频繁易姓并未让所有文书经验每次都重新发明。新朝可能改变标题和年号，却往往需要旧有的书写技术与地方知识来维持征收和诉讼。这种连续性不是政治认同的证明，而是行政生活不能中断的证明。

工匠与寺院之间也有相遇。吴越的寺塔、城市工程和堤防，南唐、后蜀的宫城与寺院，都是材料、技艺、劳力和赞助汇聚的结果。建筑遗存或题记让部分工程可见，但难以告诉我们全部劳役如何分担。因而，不能把一座塔的存在直接当作全民安定的证据；更好的读法是把它视为某些资源能够被集中、某些社群和赞助关系能够持续的迹象。

\section{河流的季节：环境不是固定的舞台}

黄河、淮河、长江、钱塘江、闽江、珠江、湘江和白藤江并不以同一种方式进入五代政治。它们在丰枯、潮汐、洪水与航行季节中改变可通行性；同一条水道能够在一时帮助运输，另一时阻断军队。国家修堤、设仓、守渡和调船的努力，正是在这种不稳定中试图建立可预测性。环境并不替任何政权作决定，却会使某些选择在某些季节变得更贵或更危险。

江淮的水网给南唐提供机动与资源，也使其防御和市场依赖河道不被对手控制。后周南征时，水的作用不仅是地图上的蓝线，而是粮船能否抵达、城池能否得到援助、部队能否在湿地行进的问题。杭州的潮汐和堤防、广州的河口和海路、白藤江的水文，则进一步说明沿水中心的能力总受地方知识限制。史料未必允许精确复原每一次水位和航程，仍足以让我们避免把战争写成无摩擦的行军。

环境史的谨慎还在于不把灾害或气候当作万能原因。收成、洪水、旱情或河道变化会影响粮食和道路，却需要通过地方制度、市场、仓储和救济才转化为政治压力。不同政权、不同城市和不同社会群体的承受力不一样。材料有时只留下朝廷关于灾异或赈济的记录，难以告诉我们实际救济到达何处；这不是可以用宏大环境叙事抹平的空白。

\section{武器之外的战争：马、舟、城与消息}

五代北方的战争常被叙述为骑兵和城战的竞争，江南与岭南的战争则容易被概括为水军行动。这样的区分有一定道理，却不能遮住军队共同面对的后勤问题。马匹需要草料、道路和适当的季节，舟师需要船材、船户、港湾和水文信息，守城需要粮仓、修缮、弓弩与居民的配合。武器只在这些条件能够持续时才有意义。后晋与契丹、后周与北汉辽的对抗之所以反复围绕城镇和关隘，正在于军事优势必须经过具体空间才能发挥。

高平之战后，后周在北方得到一段较大的行动余地，见\eventref{LT-954-111}{高平之战（954）}。这一余地不等于边防已被解决。北汉仍在太原，辽仍能影响北方通道，后周需决定把粮草、军队和官员投入何处。战后守备有时比战时冲锋更难写进史书，却更能决定胜利是否持久。若驻军无法获得稳定供给，若关口和道路不能维护，战场的结果便可能被下一次季节或下一次继承危机抵消。

水上战争也不是把陆军换成船只。南汉向交州出兵时，舰队要从广州方向经过海路和河口，见\eventref{LT-938-BACH-DANG}{白藤江之战（938）}；后周进入淮南时，粮船、渡口和沿江城镇决定部队能够怎样接近目标。水路使运输更有效，也使军队依赖某些狭窄节点。一个河口的潮汐、一段堤岸的损坏、一座城控制的码头，都能改变行动的顺序。史料很少保存完整航行记录，正不应让叙事以无阻的箭头代替这种摩擦。

消息是另一种军需。前线将领需要知道援军、粮道和敌军的位置；都城需要知道地方是否执行命令；地方官又需知道中央发生的继承与任命。驿路、使者、船只和熟悉道路的人共同构成信息网络。信息迟到并不一定造成失败，却会使行动者在不确定中作出代价高的选择。将五代的战争读作消息与供给的竞争，能够让城门内外、军队与地方社会重新出现在同一场景中。

\section{礼仪与合法性：为什么新朝需要旧语言}

五代十国的建国者很少从完全空白的政治语言开始。受禅、册封、改元、祭祀、官名和诏书都借用了此前帝国的形式。朱温受禅、王建称帝、李昪改唐、刘龑在岭南建国、孟知祥在成都称帝，分别以不同方式使用皇帝制度的语汇。旧语言并不证明新政权只是旧朝的残影；它说明行动者需要让官员、军人、邻国和地方精英能够读懂其权力要求。

合法性同样有物质的一面。改元要写入文书，任官要发出印信，祭礼要组织人员和物资，册封使节要有安全路线和礼物。若一项仪式不能被地方官、军队和市场中的人当作可执行的秩序，它的影响便有限。杨氏使用唐号，钱镠接受梁封，王延钧称帝，见\eventref{LT-907-084}{杨渥续用唐号（907）}、\eventref{LT-907-085}{吴越受封（907）}与\eventref{LT-933-MIN-EMPERORSHIP}{闽称帝（933）}，正好显示同一时代可以有不同的合法性策略。

后世史书常以统一王朝的立场评判这些策略，容易把某些称号归为正统，把另一些归为僭越。对编年叙事而言，更有用的问题是：这种语言在当时解决了谁同谁之间的关系？它是否让一支军队获得名分，让一座城市接受官员，让一个小政权避开即时冲突，或让一条贸易路线仍可使用？如此读合法性，既不放弃政治语言的力量，也不把它误认成全部社会现实。

\section{多政权的市场：交换如何穿过战争}

战争并不使交换完全停止。相反，当不同中心并存时，商人、船户、手工业者和地方市场会面对更多关津、货币、税课与风险，也可能发现新的路线和需求。开封的军队需要东南货物，江淮和两浙的城市需要北方、内地或海上输入，成都的盐货和织物需要道路，广州和福州的港口依赖沿岸交换。市场的延续并不说明民众未受战争影响；它说明许多人必须在不稳定中继续寻找生计。

货币在其中是方便而有限的媒介。后周铸钱或任何新朝的货币政策可以表明中央试图恢复交换和支付的尺度，市场中却可能同时使用旧钱、不同钱币、布帛、谷物和信用关系。出土钱币使这种多样性的一角可见，不能把一处窖藏直接解释成国家政策的成功或失败。真正的市场秩序还取决于道路、治安、价格和买卖双方的信任，而这些很少以完整统计被保存。

多政权还会改变关津的意义。一处渡口、港口或山口可能从边缘变成税收与情报的要点；一座城市可能因能够转运而获得新的地位，也可能因被封锁而迅速衰落。江陵、潭州、扬州和建州在不同阶段都可被这样理解。读者若只在国号更替中寻找变化，会错过市场如何为地方政权提供资源，又如何使地方承受额外的征收和管制。

\section{可见与不可见：给沉默留下位置}

本章反复说“不能断言”，并不是让叙事停在犹豫中。相反，材料的限度使能够确认的行动更加清楚：907年后梁、前蜀、吴越、吴和闽的不同选择；923年后唐入汴；936年石敬瑭依赖契丹；938年南汉远征受挫；945年南唐入闽；951年后周与北汉并立；954年高平；956至957年淮南战争；960年宋代建立。这些节点把多政权的年序连接起来。

不可见的是每一次征收在多少家庭中引起何种反应，每一段道路上有多少人被迫迁徙，每一座城的居民如何在新旧军队之间判断安全。没有直接材料时，把这些人写成同一种爱国者、受害者或机会主义者都不可靠。更好的方式是从他们必经的粮食、道路、市场、劳役和家族关系出发，说明为什么他们的选择会重要，同时承认选择的具体内容往往没有被保存。

这种节制也让小政权不再只是大国的配角。荆南的江陵、楚的潭州、闽的福州、南汉的广州、吴越的杭州、后蜀的成都和北汉的太原，各自拥有有限而真实的可见性。它们留下的材料有密有疏，行动空间有大有小；共同点是都不能被开封年表完全取代。编年体在这里不是中央史的另一种排版，而是让并行的城市、道路和政权在同一时间中相互照见的方法。

\section{人物不是标签：在约束中作出选择}

五代十国的人物很容易被后世整理成几种熟悉的标签：篡夺者、忠臣、割据者、降将、雄主、昏主。标签能帮助记忆，却会遮蔽行动发生的条件。朱温、李存勖、石敬瑭、刘知远、郭威、赵匡胤，分别面对军队、继承、边防和都城资源的不同组合；王建、杨氏与徐氏、钱镠、王审知、刘龑、马氏、孟氏和刘崇，也都在地区交通、精英关系和外部压力中作出选择。人物并不是被地理或制度自动推动，地理和制度却决定哪些选择可行、哪些代价会迅速显现。

石敬瑭向契丹求援尤其需要这种读法。后世从燕云问题回望936年，容易只看到一位统治者的道德失败；当时的石敬瑭却面对后唐军政危机、河东基础和取得外援的紧迫需要。援军能够帮助他建立后晋，也把新政权置入新的边防和外交关系。选择既不是没有替代方案的必然，也不是脱离军事道路和资源条件的任性。将这两点同时保留，才能理解为什么一项短期可行的决定会留下长期压力。

钱镠接受后梁册封、杨氏续用唐号，也不宜被解释为一个“忠”、一个“叛”的对照。杭州和扬州的区域条件不同，二者面对的军政集团、交通网络和外部邻国也不同。王号或唐号使各自能同不同对象沟通，背后仍要有粮食、水路、城市和地方合作来支撑。历史人物之所以值得写，不是因为他们替一种道德类型表演，而是因为他们在约束中让关系向某个方向移动，并把代价留给后来的行动者。

小政权的统治者同样如此。高季兴的江陵必须在大江通道上保持中介地位，马殷之后的楚国要处理潭州军政和家族继承，王延钧的福州要面对山海网络的分散性，刘崇的太原则须在北汉与辽的联盟中寻求安全。材料通常保存他们同中原的关系，较少保存他们同本地社会的所有互动。叙事应让他们具有行动能力，也应避免把有限材料放大成全知全能的个人策略。

\section{阅读一张不精确的地图}

五代十国的地图应画出开封、洛阳、太原、成都、江宁、扬州、杭州、福州、潭州、江陵、广州和交州等位置，以及连接它们的河流、道路和大致方向；它不应把每一年、每一县的疆界画成精确而稳定的线。政权的实际控制会随战事、册封、地方合作和运输能力伸缩，史料对不同地区的密度也不相同。一条粗略的边界若被画得过于确定，反而会让读者忘记这些变化。

地图最有用的功能不是回答“谁拥有全部土地”，而是帮助提出空间问题。为什么开封需要运河，太原离不开北方通道，江陵能成为中介，江宁和扬州需要共同考虑，杭州与广州的海路条件不同，成都的资源为何要经由山口和峡江才能进入更大世界。把事件放回这些问题中，册封、征服和联盟便不再是悬在空中的名称。

图上的箭头也应被谨慎阅读。后唐入蜀、后周入淮南、南汉向交州行动、契丹军入开封，都不是一支可以穿过空白地图的箭头。箭头所经过的是渡口、粮道、城镇、河口和地方关系；它可以标出大致方向，不能替代实际行军、船运和接收的复杂过程。正文之所以需要不断回到道路和仓场，正是为了让地图服务叙事，而不是让地图用简化的空间吞掉叙事。

\section{余波：在下一卷之前保留的问题}

本卷在960年停止，不是因为区域问题已经解决，而是因为中原中枢的转换为下一段历史打开了新的条件。宋如何面对北汉与辽，如何接收江南、吴越、南汉和后蜀，如何把军政、财政和市场重新组织进更大的体系，都是下一卷的年序。五代十国为这些问题提供的并非简单教训，而是一批仍在发挥作用的道路、城市、货币、军队和地方中介。

读者从这一章前往宋辽夏金元卷时，可以带着几个具体疑问：统一军队怎样获得长期供给；港口和运河如何进入更大财政；边地城镇是否只能由中央军政解释；地方家族、寺院、商人与书写者在新秩序中会保留什么位置；曾经依靠多重册封与联盟生存的地区将如何被接管。这样的承接使五代十国不再是统一叙事的空白地带，而成为理解下一阶段何以艰难、何以可能的必要起点。

\section{一段编年应当留下的尺度}

五代十国的读者不必在每一页记住所有年号，却应能在每个转折知道自己站在何处。907年站在开封、成都、扬州和杭州同时改变政治语言的时刻；923年站在后唐军队进入汴州而江淮、两浙仍有自身节奏的时刻；936年站在太原、燕云与契丹援军改变北方选择的时刻；938年站在广州、白藤江和交州河口之间；945年站在闽岭道路受到南唐军政压力的时刻；954至957年站在高平与淮南之间，观察后周怎样把北方胜利转化为南方战争；960年则站在陈桥之后，看见一个新中枢尚未接收完旧世界。

这种尺度防止两种相反的误读。一种把五代缩成开封皇帝的更替，仿佛南方、边地和海路只在必要时才出现；另一种把十国写成互不相干的地方故事，仿佛它们不受中原战争、册封、人口流动和市场变化影响。实际历史在两者之间：政权相互牵连，却不共享同一速度；道路和文书能够传递命令，却不能使地方马上变得相同；联盟和条约能够改变战略，却不能替代驻军、粮食和地方合作。

因此，编年叙事最后需要保留的不只是“先后”，还有“不同步”。同一年的唐号、梁号、吴越王号、蜀号、汉号或辽号，记录了行动者以不同方式组织时间；同一条河流既可以把城市联系起来，也可以在战时把它们隔开；同一份正史既能保存政权转换，也会使一些地方生活沉默。把这些不同时钟写在一起，才能使读者理解：五代十国的复杂并不是统一之前暂时的混乱，而是历史行动总在有限证据、有限资源和不均匀空间中展开的常态。
\end{causalchain}

\section{城门之后：五代十国的几次接收}

\subsection{汴州与成都：胜利必须穿过两套道路}

923年后唐军入汴州，\eventref{LT-923-092}{后唐灭后梁（923）}。城门一开，编年史可以很快改写朝代名称；真正缓慢的部分却从城门内开始。仓场里的粮食由谁封存，旧梁的官员是否继续办事，渡口的税、城内军队的口粮、来自河北和河南的文书怎样重新盖印，都不是“灭梁”二字本身能够包办的事。李存勖的军队从魏博与河东南下，靠的是已经能集结人马和传递粮秣的北方路径；进入汴州后，则要面对一座依赖黄河、汴水和东南运输联系的城市。战场上的速度，不能自动转化为日常行政的速度。

《旧五代史》的后唐纪年把入汴、任官和朝廷动作排得很清楚，《新五代史》又以不同的叙事重点整理同一轮转换。两书使我们能够确认中枢的改易，却并没有留下汴州每一坊、每一仓、每一县在数月中如何适应的连续记录。因而，读到“接收”时，不能把未被写下的社会反应补成整齐的欢呼或一律的恐惧。可以确认的是，新政权若不能让守城者、漕运者、书吏和附近军镇知道命令从何而来、俸给从何而出，入城的军事成果就会迅速变成新的不安。

两年后，后唐又进入成都，\eventref{LT-925-094}{后唐灭前蜀（925）}。这两次成功常被放在同一串“五代更替”里，仿佛只是两个都城换了主人；但汴州和成都要求的接收办法并不相同。前者处在黄河与运河相接的平原网络，后者背靠成都平原，向外的联系要经过山口、栈道与峡江。向汴州集结的部队可以依托河北、河东和中原的道路，向蜀地推进则必须计算山道的季节、关隘的守备以及粮食能否随人马进入盆地。军队抵达成都，未必表示它已取得支配蜀中所有县、盐井、乡里与商路的能力。

成都的官署、市场和寺院不因中原年号改变而停止运作；它们必须同新的征发、任命和驻军关系相接。这里不宜写成后唐将全部制度原封不动地搬进蜀地，也不宜把地方社会想成对任何来者都完全无动于衷。留下的材料主要关心朝廷如何处分王建旧部、如何宣布新秩序，较少记录普通住民怎样为军需、道路或税役重新安排一年。正是这种不对称，使“占领成都”必须同“能否在成都以外维持供应”分开理解。

后唐随后本身也陷入继承、军队和宫廷关系的紧张。它先后取得汴州和成都，并没有因此拥有一套无摩擦的跨区域财政机器。对一座都城的接收需要新旧人员合作，对一片新得区域的保持又要不断投入路粮、使者和驻军。五代的中原政权之所以看似更替迅速，并不只因为将领善于起兵，也因为每一次胜者都要继承前朝已经很紧张的运输和供给系统；一旦这些系统不能同时照顾北方边防、都城军队和远方新地，胜利留下的负担就会反过来塑造下一次危机。

\subsection{太原、燕云与草原南缘：936年的援军不是一张地图上的箭头}

936年石敬瑭在契丹援助下建立后晋，\eventref{LT-936-099}{后晋建国（936）}。后来的叙述很容易从“燕云十六州”倒推当时的每个选择，仿佛石敬瑭面对的是一张边界清楚、条件已经写好的地图。可在当时，河东起兵者首先面对的是后唐的军政压力：太原周边的军队如何聚集，援军从何处南下，谁能控制通向河北和中原的关隘，城镇能供给多久，都是先于后世道德评价的紧迫问题。契丹方面也并非抽象的“外力”；其统治者选择介入，须衡量南下通道、可获的州城、与中原新朝结盟的价值，以及部众和边地资源的承受能力。

燕云并不是一条笔直的线，而是一组山口、城堡、河谷、道路与居住其间的人。控制其中一处，意味着得到观察、屯驻或通行的条件；也意味着必须供养守军、修缮城垣、处理当地既有的官民关系。所谓割让的长期后果，正在于后晋以后的中原政权不能再把这些关口单纯视为自己的防务节点，而要把它们当作另一政治中心能够调动的前沿。这个变化并不需要把每一座州城在每一年画成确定颜色才能理解；它首先改变了双方行动时可利用的纵深和道路。

\begin{debate}{燕云的盟约能让我们确定什么}
《旧五代史》的后晋纪、《辽史》的太宗纪和《资治通鉴》都保存了936年结盟、援军和政权转换的基本年序，但它们来自不同的政治传统；盟约原件不存，十六州的具体范围与后世流传的称谓也经过整理。因而不能把后出的措辞当作谈判现场的逐字记录。不过，太原起兵依赖北方援助、契丹由此取得燕云方向的控制条件，以及这一安排改变后晋的边防处境，仍是相互参照后可以把握的历史骨架。这个限度迫使读者把注意力放在城镇、通道和联盟的后果，而非用一句传说式称呼解释全部事件。
\end{debate}

后晋建立后，联盟并未把北方变成一个安静的后方。新朝既要在开封取得官僚与财赋的承认，又要同援助者保持可以承受的关系；后者不可能只由一纸名义决定。沿边的城镇、往来的使节、军队调动和季节性的补给，都会不断检验双方的约定。后来的紧张因此并不只是个人荣辱的变化，也来自一种结构：一个依靠外援而建立的中原政权，必须在都城政治、河东旧部和北方邻国之间反复分配有限的资源与信誉。

从太原望向南方，开封并不是唯一的中心；从燕云望向南方，后晋也不是唯一能够发出命令的一方。这样的视角能避免把契丹写成只在五代皇帝更替时才短暂出现的背景。契丹拥有自己的统治重组、部众关系和北方战略，河东的军政集团也有自己的生存盘算。936年把这些行动者临时接在一起，创造了后晋，却没有消除各自目标的差异。援军帮助打开一条通向皇位的路，同时也把新王朝锁进一条需要持续维护的边防道路。

\subsection{江宁、建州与福州：南唐接手的是网络，不是一张空白的江南地图}

徐知诰于937年受禅、改名李昪，建立南唐，\eventref{LT-937-100}{南唐建立（937）}。江宁的宫城、文书和仪式让新国号获得可见的中心，但南唐接到的不是一块等待命名的领土。此前吴政权留下的将校、州县官、仓储、水运和地方关系仍在运作；淮南与江南之间的粮食、盐货和人员流动，也不会因为受禅礼成而自动改道。李昪一方需要把旧有网络纳入新的官号和任命，地方行动者则要判断这次改名会不会改变他们所面对的税役、军令与保护。

江宁与扬州之间的水路使南唐能够把政治中心同淮南相连，却也使控制依赖于堤岸、港埠、船只和地方节点。一次任命可以由都城发出，能否在遥远州县变成征收、守备或审案，则还要经过具体的人与路。南唐的国家能力不应被抽象成“江南富庶”：它既有较强的水陆资源，也要不断支付维持这些资源可用的成本。河道淤塞、驻军所需、地方将校的离合、北方政局的变化，都会使看似稳定的联系出现摩擦。

《旧五代史》《新五代史》与南唐国别材料保存了受禅和改名的线索，却大多围绕王朝中心书写。它们告诉我们杨溥让位、徐氏建国的次序，不能替江淮各地居民回答“是否接受南唐”的问题。墓志、钱币和城址偶尔让城市活动露出一角，也不能被放大成完整的财政账或舆情调查。材料的这种层次差别并不使受禅失去意义；相反，它提醒读者，受禅首先解决的是统治集团如何使命令有出处，未必立即解决命令在所有水乡、村落和市镇如何被执行。

南唐向福建推进时，这种接手的难度更加明显。945年建州失守、王延政降，\eventref{LT-945-104}{南唐取闽（945）}。从江宁的朝廷看来，闽国的内争提供了东进的时机；从建州和福州之间的山岭看来，问题却是军队、粮食、消息和任命如何穿过并不平坦的地形。建州位于山地交通的要冲，沿海港口和闽江水系又连接着另一类资源与联系。南唐能够攻下一座关键城，不表示它已用同一速度把福建各处接入江淮式的行政和供给网络。

闽末的材料尤其需要克制。后出的南唐叙事容易把内争写成胜者进入的理由，国别史的散佚又让许多地方军政集团、家族和居民的选择只留下零碎痕迹。我们可以看见政权中心的终结，却难以逐县追踪旧部何时改隶、沿海贸易如何调整、山中聚落怎样感受新的征发。把“灭闽”写作一个可确定的政治转折，并同时保留接收过程的不可见性，既避免否认南唐军事行动，也避免把福建社会写成一夜之间变成南唐的附属页。

不久后，南唐在湖南的介入又把问题推向另一套地理。潭州及湘江流域的军镇和水路，同建州的山道、福州的港口并不相同。南唐如果试图把影响延至楚地，必须在原本的江淮供给之外安排新的行军、安抚和驻守；地方势力则可以利用水路、城邑和既有关系议价。这里最重要的不是替南唐列出一条扩张清单，而是看见扩张每走一步都会带来新的行政距离。国家的边界可以在诏令中前移，真正的负担却会沿着船路、粮道和人事关系慢慢显形。

\subsection{广州与白藤江：海路把岭南和交州联系起来，也把它们分开}

917年刘龑在岭南称帝，\eventref{LT-917-090}{岭南称帝（917）}。番禺的都城地位、珠江水系和面向海上的港口条件，为南汉提供了不同于开封或太原的资源组合。岭南政权可以从海路、沿海贸易和区域农业中取得支持，却同样必须让城市、州县、军队和港口中的关系持续运作。称帝使朝廷拥有一套可以发布任命、接待使节和安排礼仪的语言；港湾、船户和沿岸市镇是否能承担军事与财赋的要求，则是另一层问题。

938年白藤江的战事把这层问题压缩在河口。南汉水军向交州方向行动，\eventref{LT-938-BACH-DANG}{白藤江之战（938）}。从广州到红河三角洲并非一条没有阻力的水上直线：舰队要经过海路，进入受潮汐、河道和地方引导影响的水口，补给也必须跟得上。即使军队抵岸，仍要面对内陆城镇和地方首领是否合作的难题。吴权一方能够凭借熟悉地形和地方组织抵抗，使一次远征没有变成稳定的接收。

越南编年叙事保存了白藤江作为建国转折的强烈记忆，汉文史书则从南汉及中原年序中简略记载。两种材料相互对读，可以确认南汉未能恢复对交州的控制，却不能把后世叙事中的每一根木桩、每一个潮位、每一项伤亡都当作精确战报。正因为战术细节不能被假装为透明，河口的空间条件反而更值得保留：水军依赖时机和向导，河口能让熟悉当地的防守者放大有限力量，远征胜败不能只用两国国号解释。

白藤江之后，岭南与交州之间没有变成没有往来的绝缘地带；贸易、人员与政治信息仍可能沿海岸移动。改变的是南汉不能把这种联系轻易转化为对交州的直接军事控制。对南汉而言，失败使广州方向的资源要重新考虑海防、港口和岭南内部的需要；对交州而言，河口的胜利为新的政治中心创造了更大的余地，却不会把地方社会的所有分歧同时消除。把这一幕放在五代十国之中，能让读者看到南方的多政权不是中原年表尽头的尾注，而有自己的海路、河口与政治时间。

\subsection{947年的空城与951年的两座都城}

947年契丹军进入开封，后晋的中枢秩序崩解；同年刘知远在太原建立后汉，\eventref{LT-947-106}{后汉建立（947）}。这两件事之间不能只用“旧朝亡、新朝兴”连接。契丹军能够抵达开封，显示北方道路、城镇和后晋军政结构已经出现决定性的裂缝；但一支军队进入都城，并不等于它已经取得维持中原日常秩序所需的一切资源。开封需要粮食、治安、官署、渡运和与周边州县的关系。远征者若不能使这些关系持续，城市的占有就会迅速暴露其限度。

刘知远的选择从河东出发。太原周边的军队、沙陀网络和北方地形，为他提供了不同于开封的基础；他不必等到所有地方都表态，便可以先以可调动的兵马和既有行政关系组织新的名义。后汉的建立说明中原秩序并没有因后晋覆亡而化为完全的空白：旧官员、军队、仓储和城市仍在，问题是谁能把它们重新编排。可这也不表示河东的军政力量天然会被所有地区接受。新朝需要离开太原的安全半径，沿着道路进入更多州县，才能把建国名号变成可持续的治理。

四年后，郭威在澶州被拥立、入开封建立后周，\eventref{LT-951-109}{后周建立（951）}；刘崇则在太原建立北汉并向辽求援，\eventref{LT-951-110}{北汉建立并求援辽（951）}。同一年出现两个相互敌对的“后汉继承”方案，最能说明政权并非只由国号决定。开封的后周继承了都城、禁军和中原官僚网络，却要处理兵变和新朝合法性；太原的北汉继承了河东兵力和刘氏关系，却必须面对后周的压力，并把对辽关系转为自身防务的一部分。

澶州与太原的空间差别很具体。澶州靠近黄河和通往开封的道路，军队拥立若要成功，必须迅速进入中枢、控制命令和城防；太原则处在北方山地与通道的交界，适合凭河东资源持守，也使北汉的安全与北方盟友的行动紧密相连。一个政权依靠都城和禁军形成中心，另一个政权依靠关隘、地方军政与外援维持边缘，两者都不是“天然正统”或“天然割据”。它们是在不同空间中回答同一个问题：当旧有继承关系断裂时，谁能供给军队、安置官员并让周围城镇承认自己的命令。

\begin{sourcenote}{后周建国不能被还原成一则兵变故事}
后周、后汉和宋代材料对郭威起兵、将士拥立与太后诏书的叙述带有各自的继统需要，《旧五代史》《新五代史》与《资治通鉴》保存的细节并不完全同调。可确定的年序是郭威取得军队支持、进入开封并建立后周；难以确定的是每位将领在何时作出何种内心决定。将叙事放回澶州到开封的行军、禁军的供给和幼主朝廷的脆弱处境，能解释事件何以可能，而不必把后出故事当作密谋的完整实录。
\end{sourcenote}

北汉求援于辽也不该被轻率地写成一方完全失去主动性。刘崇和河东集团需要面对后周的军事压力，辽则有借太原影响中原北部的战略利益；联盟给双方带来可用的资源，也给双方带来相互牵制。对太原附近的城镇而言，联盟意味着驻军、道路与供给可能重新被安排；对辽而言，援助北汉并不等于能直接管理河东每一处地方。正是在这种不完全控制中，北汉得以存续，后周也无法把北方问题当作一次建国后即可清除的余波。

\subsection{高平之后：一场战役获得的不是永久安全}

954年，高平成为后周、北汉与辽竞争的会合点，\eventref{LT-954-111}{高平之战（954）}。新即位的柴荣必须面对一个非常实际的局面：若北汉与辽的压力穿过河东南下，后周在开封的立足会立刻受考验；若禁军中的将领观望，刚建立的朝廷就难以把军令延伸到前线。高平的胜利因此为后周赢得了行动余地，但这份余地的内容不是“北方从此安定”，而是让朝廷有可能重新安排军队、奖惩和边防资源。

战场叙事最容易把复杂的过程压成几个勇将和一次临阵转折。《通鉴》和人物传记确实留下了有关退却、督战和胜败的鲜明片段，其他史书对辽军投入和战术过程的描写又有差异。兵数、阵形和伤亡缺乏可以一一核实的共同底本，不能为了制造壮阔场面而写成整齐数字。更可靠的读法是注意高平所在的北方道路：军队要在关隘、河谷与城镇之间移动，马匹和士卒要得到补给，胜者战后还要把驻守、巡防和奖惩落在可执行的制度上。

柴荣对后周军队的整顿，因而既有即时的军事目的，也有财政和社会的成本。强化军纪可以减少临战瓦解的风险，却需要一套能支付军饷、补充装备和处理违令者的中枢；把兵力集中于北方，可能又会压缩其他地区的守备和地方的余粮。史书常把改革写成君主的果断，实际的效果则要经过将领、军吏、州县和供应者共同实现。没有这些中介，命令只能停留在都城或军营。

对北汉而言，高平失利没有使太原从地图上消失。刘崇政权仍可依托河东地形、城镇和对辽关系维持抵抗；对辽而言，一次战场挫折也不等于放弃南方通道。后周必须继续在北方投入注意力，这正是柴荣后来能够转向淮南时仍须计算的约束。战役改变了各方下一步的选择，不会替各方解决粮道、联盟和继承的长期问题。

\subsection{寿州、淮河与长江：南征把两种国家能力拉到同一条水路上}

956年后周军攻入淮南，\eventref{LT-956-113}{后周攻入淮南（956）}。从开封看去，南征似乎是高平后军力增长的自然延续；沿着淮河行军，才会发现它是对另一种国家能力的检验。北方军队需要把粮食、命令、渡河器具和增援送进水网密布的地区，不能仅凭此前在河东的经验行事。南唐则要依靠寿州、濠州等城防、水道和当地供给拖慢对手，同时设法让江宁的决策能够传到前线。双方的优势都不是抽象的“强”或“富”，而要在船只、堤岸、渡口、仓储和守军的协作中才能出现。

寿州之类的城池所以反复进入史书，不只是因为它们被攻守，更因为它们把淮河两岸的路线收束到可争夺的节点。围城的一方必须保证后方粮道未被切断，守城的一方则要衡量粮食、援军和突围的可能；城外的居民、船户和附近乡里也会面对新的征发、运输和避难压力。两《五代史》和《资治通鉴》能够较清楚地勾勒战役进程，但对于淮南家庭实际承担了多少劳役、多少人迁离或多少田地荒废，保存的材料远不连续。叙事可以指出战争必然会触及这些生活条件，却不能擅自替他们填上统一的数字和感受。

南唐的江北防线并非单纯一道向北的墙。它同时连接着江淮的市场、盐粮运输、州城官署和长江方向的退路。后周军若取得某些节点，南唐失去的不只是地图上的地名，而是使北部资源与江宁中心相互流动的条件。反过来，后周越向南推进，自己的补给线也越长，越依赖新占州县能否提供渡运、驻军和信息。战争把两个政权的行政能力放在同一片湿润、易变的空间中比较，而不是让某一方以单一优势压倒另一方。

957年，南唐向后周称臣并割让江北诸州，\eventref{LT-957-114}{南唐割江北（957）}。外交文书可以给这轮战争一个收束，接收却仍在此后进行。后周必须决定哪些旧官可以留用、何处设防、税粮怎样转输；南唐也要重估失去江北后的边防、财政和政治信誉。对原在南唐统辖下的州县而言，最直接的变化未必首先是“我是何国之民”，而可能是新来的军队住在哪里、谁征粮、渡口向谁纳税、旧契约和官印是否还有效。称臣的语汇属于朝廷之间的关系，不能被误读为一份地方认同的总表。

这场战争也改变了吴越等邻近政权观察中原与江南的方式。杭州的统治者需要考虑长江下游和海路的安全，商人与使节要判断哪些路线仍可通行，南唐则不能再把北部缓冲地带视为理所当然。吴越并未因后周南征而消失于场景之外；恰恰相反，区域力量的重新排序使它在外交、港口和沿海防务上有新的判断。多政权的历史从不只发生在交战双方，旁观者也要为道路、贸易和册封关系的变化付出选择的成本。

\subsection{幼主、北伐与陈桥：960年结束的是后周，不是并行世界}

959年柴荣去世，幼主继位，\eventanchor{LT-959-117}{周世宗去世与后周幼主继位（959）}。前数年积累的军队整顿、北方防务和淮南接收尚未变成一种不依赖具体统合者的自动秩序。开封朝廷要在辅政大臣、禁军将领和各地军政之间维持平衡；北方仍有北汉与辽，南方则有南唐、吴越、南汉、后蜀等各自的局面。幼主并不必然导致政权覆亡，但继承把原先能够暂时压住的军队归属和信息不对称重新暴露出来。

陈桥兵变发生在这样的条件下。960年赵匡胤被将士拥立、进入开封受禅建立宋，\eventref{LT-960-118}{陈桥与宋代建立（960）}。后世最熟悉的是黄袍加身的画面；这一画面来自宋代建国后反复整理的叙述，不能当作我们能够旁观的现场。更可把握的是军队因北方警报而集结、禁军主帅拥有关键位置、后周幼主朝廷难以完全控制这一过程，以及赵匡胤最后取得开封中枢。兵变能够转成受禅，正因为它不只发生在军营，还需要都城城防、官僚文书与旧有行政网络允许权力交接不致立刻陷入全面冲突。

宋的建立并没有在同一天抵达所有区域。太原的北汉仍然存在，南唐仍以江宁为中心，吴越、南汉、后蜀及其他地方政权仍要处理自己的军队、港口、税源和外部关系。对这些政权而言，陈桥首先是北方中枢换了一个可以谈判或防备的对象，而不是自身历史的句号。对开封之外的商人、僧侣、船户、地方官和农户而言，变化会沿着命令、征发、道路和市场逐步传来，其速度不必同正史的改朝日期一致。

五代十国在960年留下的，因而不仅是下一王朝的征服目标，还有一套已经形成的区域现实：燕云和河东的北方压力，淮河与长江之间的接收难题，江南的水路财政，岭南与交州的海河联系，福建山海网络，以及成都、江陵、潭州、杭州等城市各自的中介位置。宋初的统一将要面对这些现实，而不能把它们从头创造出来。把本卷停在陈桥，是为了让读者看清一个新中枢出现的时刻；不把这一刻写成万事归一，才是对五代十国并行年序最基本的尊重。

\subsection{杭州的堤岸与册封文书：吴越不在中原战争之外，也不由中原战争决定}

907年钱镠接受后梁所授吴越王号，\eventref{LT-907-085}{吴越受封（907）}。一纸册封把杭州同开封新立的后梁连接起来，但杭州的政治位置并不从属于这段文辞的单一方向。两浙的水网、钱塘江口和沿海航路，使这里既要面对内陆的军政变化，也要管理海潮、港埠、粮货和沿岸安全。接受王号可以让吴越与北方政权建立一条可供使节、任命和贸易使用的名义通道；它不能替代地方官、军队、船户和市镇在日常中实际维持的秩序。

杭州不是一座只在册封时才出现的城市。江潮和河流决定堤岸、渡口与航行的条件，城市与腹地之间的粮食、木材、盐货和人员往来，也要经由具体的水陆节点。地方政权若要在这里长久存在，便不能只把资源用于宫廷礼仪或外部战争；它还须让运输可用、城防有人承担、市场有可以预期的规则。材料没有留下完整的年度收支表，不能把吴越的“经营”写成每一项工程都由中央精确规划；但城市遗存、钱币、碑刻及后出的国史线索，足以提示我们：水利、港口与地方行政不是背景，而是政权能够同外界议价的条件。

吴越的选择也不能被浪漫化为一种天然的“偏安”。北方改朝会影响册封语言和陆路往来，江淮局势会改变杭州与扬州、江宁之间的风险，沿海交换又会把本地政策同更远的港口相连。钱镠及其继承者需要判断何时维持旧有名义、何时调整对邻国的关系；地方社会则要在这些调整中承担防务、运输和征收的具体后果。吴越能够保有相对长的延续性，不等于其内部没有压力，而是说明它把城市、水网和外交关系维持成了一种可以反复调节的区域能力。

《吴越备史》及相关传记保存了王室与朝廷往来的丰富材料，却多在后来的政治记忆中成形；中原正史又习惯以册封和朝贡观察吴越。两类材料都不宜被当作杭州所有人的声音。它们能让我们看见政治中心如何说明自己的位置，却难以逐村说明税役如何落下、沿海居民怎样看待王号。把材料的边界保留在这里，并不是缩小吴越的历史，而是防止一份朝廷文书把一整片水乡压成无声的背景。

\subsection{江陵的码头：荆南以中介位置而不是疆域大小进入年序}

江陵位于长江中游的要冲。后唐在924年册封高季兴为南平王，\eventref{LT-924-093}{高季兴受封南平王（924）}。对北方朝廷来说，这一安排可以避免立刻为控制江陵而投入一场代价高昂的战争；对高季兴来说，王号为本地军政集团提供了一种可以同邻国沟通的名义，却不要求他把城内的仓储、船只和兵员完全交给后唐。册封因而不是中央向边缘单向施予的静态事实，而是双方在交通节点上暂时达成的分工。

长江在此处把不同方向接在一起。向东可以通往江淮水路，向西要考虑峡江与蜀地，向北又有荆襄陆路和通向中原的路线。江陵的价值不只在城墙以内，而在于船只、渡头、粮食和消息是否能在这些方向之间被调度。小政权的空间往往比地图上的国界更具弹性：它也许不能广泛占有土地，却能影响谁可以通过、谁能补给、谁需要协商。高季兴及其后继者的困难同样来自这一点。通道一旦因战争、封锁或邻国关系变化而受阻，作为中介所得到的收入和安全就会同时受压。

从江陵居民的角度看，所谓“依附”或“独立”不能替代更具体的问题。哪个使节可以入城，哪支军队要过江，谁征收仓场与渡口的费用，城防需要多少劳役，都会改变城市生活。现有史书更善于记录高季兴同后唐、后晋或南唐之间的往来，较少保存码头工人、船户和乡村供给者的自述。我们不应假设所有人都因中介地位获益；但也不该把江陵写成被大国随意穿越的空地。它的政治作用恰恰来自有人能够在这座城中组织通行、供给与防卫。

荆南的存在使五代十国的时间不再只由“谁占开封”来分段。后唐、后晋、后汉、后周的年号在北方更替时，江陵必须一次次判断新的中枢是否仍值得保持关系；南唐、楚、吴越和蜀地的变化又会让长江上的选择重新计算。对小政权来说，持续生存不是远离所有冲突，而是把不能避免的通道压力转化为可以谈判的地位。这一能力有其成本，也有其脆弱处，正因如此，荆南值得在五代年序中拥有自己的场景。

\subsection{潭州与湖南水系：楚亡之后，南唐也没有得到一块可以直接调用的腹地}

楚国以潭州为中心，连接湘江、洞庭湖及通向岭南、江淮的水陆路线。它的政治资源来自军镇、地方精英、农产与水上运输的结合，而不是从开封年表自然延伸出来的一条支线。马氏统治下的湖南在战争时代能够组织区域秩序，正因为潭州并非孤城：粮食、木材、船只、盐货和人员可以沿水系来去，地方军政也能借此形成自己的重心。可这种联系并不总是优势。水路方便调运，也使要道成为争夺、征收和封锁的对象。

951年前后楚国内部的继承和军政矛盾加深，南唐在952年介入湖南，\eventref{LT-952-SOUTHERN-TANG-CHU}{南唐介入湖南（952）}。从江宁看来，这似乎是把影响力推入一片有河流相连的地区；从潭州及其周围看，却是一项需要重新处理地方军队、仓储与水路的接收。南唐的军队可以沿着某些通道抵达关键城市，却仍需知道哪一段河道适宜船运，哪一处旧部可以安置，哪些地方关系会因新的任命而反弹。楚的政权中心消失，并不意味着湖南马上变成一块同江宁协调一致的资源库。

这件事同945年入闽形成鲜明对照。福建的难处在于山岭与港口并存，湖南的难处则在于江湖水系把城市、乡里和军事路线连得很紧，却也让控制者必须持续掌握船运和节点。两地都让南唐看到扩张的边界：军队可以利用对手内争进入，但新政权能否把占领转成稳定的税源和守备，要由地方环境和人际关系决定。用同一句“南唐兼并”概括两件事，会遮住建州、福州、潭州和湘江各自不同的政治地理。

南唐对湖南的影响后来又受到中原、淮南战争和内部资源分配的限制。一个政权同时试图管理江淮、福建与湖南，就必须让船队、军队、官员和财政在很长的距离中保持联系；每一处新得地区都可能要求额外的安抚和给养。史料很少让我们精确计算这种成本，却反复显示了一个基本事实：多政权时代的“扩张”并不是把地图着色，而是把已有的能力投入一个对自己并不完全熟悉的地方，并承受因此产生的新的不确定。

\subsection{成都再建国：后蜀的稳定首先是一种山地与平原之间的日常工作}

934年孟知祥在成都建后蜀，\eventref{LT-934-098}{后蜀建立（934）}。前蜀覆亡不到十年，蜀地又出现新的政权中心，这种紧凑的年序很容易诱使读者把成都想成可以不断被不同王朝随手取得的舞台。事实上，前蜀、后唐和后蜀所面对的成都平原、山口和峡江并没有消失；每一方都必须同这些条件重新建立联系。孟知祥能够在蜀地取得位置，不只因为中央权力出现缝隙，也因为地方军政、道路和资源已经形成了不完全由洛阳或开封控制的组合。

成都平原的农业与城市供给给政权提供基础，周围山地却使向外扩张、接受中央命令或应付外来军队都要经过有限通道。山道能够保护盆地，也会使外部物资、使者和增援难以及时进入；峡江可把蜀地接到长江网络，也可能在局势紧张时变成昂贵而危险的路线。后蜀的统治者要维持的不只是宫廷中的名号，还包括城内外的粮食转运、军队驻守、同地方精英的关系和对山口的防护。没有这些日常工作，“蜀中富庶”便只是后人的空泛印象。

后蜀与中原的关系也不应被设想成彻底隔绝。使节、商货、僧侣、逃亡者和消息仍可能沿道路与水路流动，北方政局会影响蜀地对边防与外交的判断。与此同时，成都的选择还会影响江陵、荆襄和长江上游的力量平衡。后蜀并非只是在等待960年以后的统一战争，而是在本卷所涵盖的每一年中经营自身的防务和社会关系。材料的局限使我们不能把王室叙事扩展成全蜀百姓的共同经验；但将成都放回平原、山道和峡江之间，已经足以改变“十国只是中原乱世余波”的误读。

把这些场景并置，五代十国才显出它本来的尺度。汴州的仓场、太原的关隘、杭州的江潮、江陵的码头、潭州的江湖、成都的山口、广州外海与白藤江河口，不是为一条中原帝王年表添上的地名。它们各自决定了军队能走多远，粮食能在哪里停留，命令会被谁接住，外来的王号和新旧国号又会在何处遇到抵抗、协商或沉默。不同政权有不同的资源，也有不同的资料保存条件；让它们在同一年代并列，不是为了制造一张热闹的名单，而是为了重新看见政治选择总要落在道路、城市和实际生活上。

这也是阅读五代十国时应有的一种节奏。读到后唐入汴，不要立刻跳到宋的建立，而应停在汴州的官署和仓场，追问接收怎样完成；读到石敬瑭求援，应回到太原、燕云和契丹南下的道路，追问短期同盟如何改变每一方的防务；读到南唐进入建州或潭州，应沿着山路、河港与旧军镇观察新朝的命令能走到哪里；读到南汉军在白藤江受挫，应记住河口的潮汐与地方组织能把远征舰队挡在政治控制之外。如此，后周的高平、淮南战争和陈桥受禅也不再是通向统一的几枚台阶，而是不同政权在有限资源中连续作出的选择。史书保存的常是皇帝、将领与使节的名字，叙事却应让读者同时看见支撑或限制这些名字的船、粮、城、路和尚未留下姓名的人。

\subsection{接收江北时，后周接收的不是一条已经服从的河岸}
936年石敬瑭借契丹之力建后晋并承诺燕云，\eventref{LT-936-099}{燕云安排}的原始文书未存，后出中原叙事带有强烈道德判断。可确认的是，太原、河北与北方关隘的军政竞争需要外援，而外援改变了谁能经过这些关口。946年契丹入开封、后晋亡，\eventref{LT-946-105}{开封易手}，也说明占领都城不等于能长久供给和统治中原。城门的得失必须同仓场、退路和地方守军的选择一起读。

951年后周建立，954年高平获胜，\eventref{LT-951-109}{郭威建周}和\eventref{LT-954-111}{高平之战}提供了重建北方军事的机会。战役阵形和伤亡数被后来的传记戏剧化，胜利之后的粮饷、马匹和将校安置却不会自动完成。955年整顿寺院、铸周元通宝，\eventanchor{LT-955-112}{后周废并寺院（955）}与\eventanchor{LT-955-LATER-ZHOU-COINAGE}{铸周元通宝（955）}都是寻找可支配资源的办法；诏令与窖藏都不能证明全国效果，代价仍由具体寺院、工匠、市场和地方社会分别承担。

956至958年，后周南下，南唐称臣割江北州，\eventref{LT-957-114}{南唐割地}与\eventanchor{LT-958-115}{后周接收江北（958）}不能被压成地图上的一笔。扬州等地需要重设军政和漕运据点，居民要面对新的催粮者、守军和渡口检查。959年北伐因病收束，960年陈桥受禅，\eventref{LT-959-116}{短期北伐}与\eventref{LT-960-118}{宋建立}给中原年序写下终点，却不使杭州、成都、广州、潭州和太原的区域道路同时消失。五代十国的遗产正是这些不同政权曾经组织过的财政、城市和社会联系。
